\documentclass[11pt]{article}
\usepackage[a4paper,margin=25mm]{geometry}
\usepackage{amsmath,amssymb,amsthm,mathtools,array}
\usepackage[hidelinks]{hyperref}
\allowdisplaybreaks[2]
\setlength{\emergencystretch}{2em}
\title{Consecutive log-concavity violations in independence sequences of trees}
\author{Leo Gao\\[-1pt]{\small AI-assisted research}}
\date{November 26, 2026 \quad$\cdot$\quad Research draft}
\begin{document}
\maketitle
\begin{abstract}
For every positive integer $L$, we construct infinitely many trees of diameter seven whose independence sequences contain $L$ consecutive strict violations of log-concavity. An explicit tree on 14,662 vertices exhibits eight consecutive violations, verified by exact computation.
\end{abstract}
\section{Background}

For a finite graph \(G\), its independence polynomial is

\[
I(G;x)=\sum_{j=0}^{\alpha(G)}i_j(G)x^j,
\]

where \(i_j(G)\) counts independent vertex sets of cardinality \(j\), and \(\alpha(G)\) is the independence number.

The coefficient sequence is log-concave if

\[
i_j^2\geq i_{j-1}i_{j+1}
\]

at every interior index.

Previous results include the following.

\begin{center}
\small
\begin{tabular}{@{}>{\raggedright\arraybackslash}p{.17\linewidth}>{\raggedright\arraybackslash}p{.78\linewidth}@{}}
Year & Development \\
\hline
1987 & Alavi, Malde, Schwenk and Erdős conjecture unimodality for trees. \\
2023 & Counterexamples to the stronger log-concavity conjecture are identified~\cite{KL}. \\
2025 & David Galvin establishes violations arbitrarily far from the terminal coefficient~\cite{Galvin}. \\
2025 & César Bautista-Ramos establishes arbitrarily many violations, without requiring consecutiveness~\cite{BR}. \\
June 2026 & Bautista-Ramos, Guillén-Galván and Gómez-Salgado establish examples with up to five consecutive violations~\cite{BGG}. \\

\end{tabular}
\end{center}

\section{Main theorem}

\textbf{Theorem.} For every positive integer \(L\), there exist infinitely many pairwise non-isomorphic trees \(T\) of diameter seven such that, for some integer \(r\),

\[
\boxed{
i_j(T)^2<i_{j-1}(T)i_{j+1}(T)
}
\]

simultaneously for every

\[
j=r,r+1,\ldots,r+L-1.
\]

Consequently, no universal upper bound exists on the length of a consecutive run of log-concavity violations in the independence sequences of trees.

\section{The construction}

Define \(T(m,n,k,h)\) using two adjacent central vertices, \(u\) and \(v\).

\begin{itemize}
\item Attach \(m\) subdivided \(k\)-edge stars to \(u\), joining their centers to \(u\).
\item Attach \(n\) identical subdivided stars to \(v\).
\item Attach \(h\) additional paths of length two to \(u\).
\item Attach two additional leaves to \(v\).

\end{itemize}
Every subdivided star consists of a center and \(k\) paths of length two emanating from it.

The additional paths at $u$ control the relative contributions of the cases in which $u$ is selected and in which neither central vertex is selected.

The total number of vertices is

\[
|V(T)|=(m+n)(2k+1)+2h+4.
\]

For \(m,n\geq 2\), its independence number is

\[
\alpha=(m+n)(k+1)+h+2.
\]

Its diameter is exactly seven.

\section{Proof}

\subsection{Lemma 1 --- Exact polynomial identity}

Define the reversed independence polynomial

\[
R_T(y)=y^\alpha I(T;1/y)
\]

and let

\[
H_k(y)=(1+y)^k+y(2+y)^k.
\]

Partition independent sets according to whether they contain \(u\), contain \(v\), or contain neither. Since \(u\) and \(v\) are adjacent, these are the only possibilities.

This gives the exact identity

\[
\begin{aligned}
R_T(y)={}&H_k(y)^{m+n}(2+y)^h(1+y)^2\\
&+y^{m-1}H_k(y)^n
(2+y)^{km}(1+y)^{h+2}\\
&+y^{n+1}H_k(y)^m(2+y)^{kn+h}.
\end{aligned}
\tag{1}
\]

To verify the identity, a subdivided star has ordinary independence polynomial

\[
(1+2x)^k+x(1+x)^k.
\]

If its center is forbidden, the polynomial becomes \((1+2x)^k\).

The three cases in (1) have maximum independent-set sizes \(\alpha\), \(\alpha-m+1\), and \(\alpha-n-1\), respectively. Reversing their polynomials therefore produces the stated powers of \(y\).

This also verifies the claimed independence number. \(\square\)

\subsection{Lemma 2 --- A limiting polynomial family}

Fix integers \(m,n\geq2\) and an integer \(\ell\). Set

\[
h=k-\ell.
\]

For sufficiently large \(k\), this is nonnegative.

Define the rescaled polynomial

\[
Q_k(y)=2^{-h}R_T(2^{-k}y).
\]

Then, coefficientwise as \(k\to\infty\),

\[
\boxed{
Q_k(y)\longrightarrow
(1+y)^{m+n}
+2^\ell y^{m-1}(1+y)^n.
}
\tag{2}
\]

\textbf{Proof.} First observe that

\[
\begin{aligned}
H_k(2^{-k}y)
={}&(1+2^{-k}y)^k\\
&+y(1+2^{-k-1}y)^k
\end{aligned}
\]

converges coefficientwise to \(1+y\).

Similarly, the other normalized binomial factors in (1) converge coefficientwise to 1.

The multiplicative scalar in the second term becomes

\[
2^{-k(m-1)}2^{km-h}
=2^\ell,
\]

while the scalar in the third term becomes

\[
2^{-k(n+1)}2^{kn}=2^{-k},
\]

which converges to zero.

Substitution into (1) proves (2). \(\square\)

\subsection{Lemma 3 --- The limiting family contains arbitrarily long strictly log-convex intervals}

Consider

\[
F_{M,\ell}(y)
=(1+y)^{2M}+2^\ell y^{M-1}(1+y)^M.
\]

Take \(M\) even and set

\[
J=\frac{3M}{2}.
\]

Write its coefficients as

\[
f_j=A_j+2^\ell B_j,
\]

where

\[
A_j=\binom{2M}{j},
\qquad
B_j=\binom{M}{j-M+1}.
\]

Choose

\[
\ell=
\left\lfloor
\log_2\frac{A_J}{B_J}
\right\rfloor.
\]

Consequently,

\[
t_M=\frac{2^\ell B_J}{A_J}
\in(1/2,1].
\tag{3}
\]

For every fixed integer offset \(s\), the standard ratios of consecutive binomial coefficients give

\[
\frac{A_{J+s}}{A_J}\longrightarrow 3^{-s},
\]

and

\[
\frac{B_{J+s}}{B_J}\longrightarrow 1.
\]

Thus,

\[
\frac{f_{J+s}}{A_J}
=3^{-s}+t_M+o(1).
\tag{4}
\]

The error tends uniformly to zero over any fixed finite collection of offsets.

Now define

\[
g_s=3^{-s}+t,
\qquad t\in[1/2,1].
\]

A direct calculation gives

\[
\begin{aligned}
g_{s-1}g_{s+1}-g_s^2
&=\left(3+\frac13-2\right)t3^{-s}\\
&=\frac43t3^{-s}\\
&>0.
\end{aligned}
\tag{5}
\]

For any prescribed finite interval of offsets, this expression has a strictly positive lower bound uniform in \(t\in[1/2,1]\).

Equations (3)--(5) imply that, for every positive integer \(L\) and all sufficiently large even \(M\),

\[
f_{J+s-1}f_{J+s+1}>f_{J+s}^2
\]

for all \(s=0,\ldots,L-1\). \(\square\)

\subsection{Completion of the theorem}

Given \(L\), choose \(M\) and \(\ell\) as in Lemma 3.

Apply Lemma 2 with \(m=n=M\). Because the required inequalities are strict and depend continuously on finitely many coefficients, they persist in \(Q_k\) for every sufficiently large \(k\).

Finally, log-concavity and its strict violation are invariant under geometric coefficient rescaling:

\[
b_j=ca^j r_j
\]

implies

\[
b_{j-1}b_{j+1}-b_j^2
=c^2a^{2j}(r_{j-1}r_{j+1}-r_j^2).
\]

Reversal of the coefficients also preserves consecutive violations, merely reversing their positions.

Therefore the original independence polynomial \(I(T;x)\) has at least \(L\) consecutive log-concavity violations.

Increasing \(k\) produces trees of strictly increasing order, giving infinitely many pairwise non-isomorphic examples, all of diameter seven.

\[
\boxed{\text{Q.E.D.}}
\]

\section{An explicit tree}

We tested finite realizations of the limiting polynomial family using exact integer arithmetic.

A successful parameter choice is

\[
\boxed{
m=46,\quad n=96,\quad k=51,\quad h=16.
}
\]

The tree has 14,662 vertices, 14,661 edges, and independence number 7,402. The run consists of eight violations.

Define the violation ratio

\[
V_j=\frac{i_{j-1}i_{j+1}}{i_j^2}.
\]

Every value strictly above 1 constitutes a violation.

\section{Measured coefficient curvature}

\begin{center}
\small
\begin{tabular}{ll}
Coefficient index \(j\) & Violation ratio \(V_j\) \\
\hline
7,304 & 0.983700096 \\
\textbf{7,305} & \textbf{1.000070050} \\
\textbf{7,306} & \textbf{1.022410490} \\
\textbf{7,307} & \textbf{1.046287768} \\
\textbf{7,308} & \textbf{1.061621514} \\
\textbf{7,309} & \textbf{1.059551860} \\
\textbf{7,310} & \textbf{1.042087410} \\
\textbf{7,311} & \textbf{1.019666129} \\
\textbf{7,312} & \textbf{1.000749025} \\
7,313 & 0.987902454 \\
\end{tabular}
\end{center}

The displayed ratios are rounded for readability. All eight inequalities were verified using exact integers, without floating-point comparisons.

An independent computation constructed the graph in memory and applied a rooted-tree dynamic programming algorithm. Its independence number and the 99 highest-degree coefficients agreed exactly with the explicit polynomial calculation.

\section{Reproduction}
The accompanying \texttt{verify.py} uses exact integer arithmetic to check the eight strict inequalities for this tree.

\section*{Disclosure}
This manuscript was prepared with AI assistance. The arguments and priority claims are open to independent review.

\begin{thebibliography}{9}
\bibitem{KL} A. Kadrawi and V. E. Levit, The independence polynomial of trees is not always log-concave starting from order 26, arXiv:2305.01784 (2023). \href{https://arxiv.org/abs/2305.01784}{Source.}
\bibitem{Galvin} D. Galvin, Trees with non log-concave independent set sequences, arXiv:2502.10654 (2025). \href{https://arxiv.org/abs/2502.10654}{Source.}
\bibitem{BR} C. Bautista-Ramos, Multiple breaks of log-concavity in the independence polynomials of trees, arXiv:2511.00334 (2025). \href{https://arxiv.org/abs/2511.00334}{Source.}
\bibitem{BGG} C. Bautista-Ramos, C. Guillén-Galván and P. Gómez-Salgado, Linear recurrences for non-log-concave independence polynomials of trees, Graphs and Combinatorics (2026). \href{https://doi.org/10.1007/s00373-026-03054-4}{Source.}
\end{thebibliography}
\end{document}
